\section{Additional Results}
\label{app:extra}




\paragraph{Solver lifecycle.}
\cref{tab:lifecycle} decomposes a re-solve on the HUNL turn subgame and on Libratus
endgame~4 into every phase a caller pays, from reading the specification to the
in-place \texttt{reset()}, \texttt{set\_payoffs} and \texttt{set\_root\_ranges}
calls that precede each solve (\cref{app:compiler}), and reports the measured
wall-clock of 1, 2, 10 and 100 back-to-back solves that reuse one solver against the
rebuild-per-solve model estimate.
\input{figures/TABLE_T14_lifecycle.tex}

\paragraph{Peak GPU memory.}
\cref{tab:memory} lists peak allocated and reserved GPU memory per game for the
1{,}000-iteration CFR$^{+}$ solve behind \cref{tab:matrix}.
\input{figures/TABLE_T4_memory.tex}

\paragraph{Timing dispersion.}
\cref{tab:spread} reports the repetition counts and the spread each \cref{tab:matrix}
median was taken over. Repetitions are independent processes, and both spread columns
are independent maxima across games, relative to each game's median: the interquartile
range and the full $\max-\min$ range. Parentheses name the range-maximizing game.
Spread concentrates on the smallest trees, where an iteration is short enough for
process-level noise to matter, and on the eager path, whose per-operation dispatch stays
sensitive to host scheduling at every size. On the four largest games every other arm
holds within $10\%$, and the graph path within $2\%$.
\input{figures/SNIPPET_exclusions.tex}

\input{figures/TABLE_T2b_spread.tex}

\paragraph{Kim (2026): precision.}
Double precision behaves as it does on GPU-CFR (\cref{sec:ladder}), for the same reason: the
iteration is launch bound. On \texttt{leduc\_poker} and
\texttt{liars\_dice}, float64 costs $0.999\times$ and $0.991\times$ the float32 time
and changes exploitability by at most $6.3\%$
($6.636\mathrm{e}{-3}$ vs.\
$6.240\mathrm{e}{-3}$ and $3.441\mathrm{e}{-3}$ vs.\
$3.445\mathrm{e}{-3}$).

The same launch pattern still performs well relative to OpenSpiel on larger games:
Kim's GPU backend is $11$--$176\times$ faster on the four larger shared games. On
Kuhn and Dark Hex, fixed launch overhead dominates both GPU implementations.

\subsection{Cross-implementation check of the update-rule study}
\label{sec:convergence}

The cross-implementation grid in \cref{tab:crossframework} checks the ranking pattern of
\cref{tab:convergence} (\cref{sec:results}).
Across the 10 comparisons with a unique winner on both sides, the winner agrees in 10;
2 further comparisons are excluded for a tie on at least one side. Because averaging
conventions vary across implementations, this comparison concerns within-backend rule
rankings. We adjudicate close rankings on the CPU path, whose reduction order is fixed,
and use the GPU path for wall-clock timing; \cref{app:correctness} reports the
repeated-run GPU analysis. Values are ranked only within a backend, game, and update
order; averaging tags make explicit why absolute values are not compared across
backends.

\input{figures/TABLE_T10_crossframework.tex}

\paragraph{LiteEFG: both arms.}
LiteEFG ships four update rules: vanilla CFR, CFR$^{+}$, DCFR and PCFR$^{+}$. Only
its vanilla CFR uses simultaneous updates; the other three alternate. None of its arms
matches GPU-CFR on both axes at once, since the latter is CFR$^{+}$ with simultaneous
updates, so we measured the two closest ($n=10$ per cell). \cref{tab:matrix} reports the
simultaneous vanilla arm, which yields $185\times$ and $213\times$ on Liar's Dice and
Battleship. The alternating CFR$^{+}$ preset costs $1.28$--$1.47\times$ more per
iteration and therefore yields $265\times$ and $306\times$ on those two games. The preset reaches lower at matched iteration counts (for instance
$1.24\mathrm{e}{-4}$ vs.\ $1.85\mathrm{e}{-2}$ on Liar's Dice), but that gain mixes
the variant with the update order and so is not attributable to either alone; we
therefore compare per-iteration cost at matched update order and leave the
convergence comparison to \cref{fig:curves}, which plots exploitability against
wall-clock time for the vanilla CFR arm.

\paragraph{Measurement provenance.}
Every timing number in the paper comes from the same gated campaign: each
measurement cell ran as its own process with its GPU and pinned cores isolated from the
host's other tenants, with
the load average sampled before and after and the cell rejected if it exceeded a
fixed threshold (\cref{sec:setup}). Reported values are medians over 20
repetitions for GPU-CFR and the Ours CPU arm and 10 for Kim's GPU path and the CPU
baselines, with ranges in the tables.
